\documentclass[aps,preprint]{revtex4}%
\usepackage{amssymb}
\usepackage{amsfonts}
\usepackage{amsmath}
\usepackage{graphicx}%
\setcounter{MaxMatrixCols}{30}
%TCIDATA{OutputFilter=latex2.dll}
%TCIDATA{Version=5.50.0.2960}
%TCIDATA{CSTFile=revtex4.cst}
%TCIDATA{Created=Friday, August 14, 2015 10:56:37}
%TCIDATA{LastRevised=Thursday, October 22, 2015 15:16:13}
%TCIDATA{<META NAME="GraphicsSave" CONTENT="32">}
%TCIDATA{<META NAME="SaveForMode" CONTENT="1">}
%TCIDATA{BibliographyScheme=Manual}
%TCIDATA{<META NAME="DocumentShell" CONTENT="Articles\SW\REVTeX 4">}
%BeginMSIPreambleData
\providecommand{\U}[1]{\protect\rule{.1in}{.1in}}
%EndMSIPreambleData
\newtheorem{theorem}{Theorem}
\newtheorem{acknowledgement}[theorem]{Acknowledgement}
\newtheorem{algorithm}[theorem]{Algorithm}
\newtheorem{axiom}[theorem]{Axiom}
\newtheorem{claim}[theorem]{Claim}
\newtheorem{conclusion}[theorem]{Conclusion}
\newtheorem{condition}[theorem]{Condition}
\newtheorem{conjecture}[theorem]{Conjecture}
\newtheorem{corollary}[theorem]{Corollary}
\newtheorem{criterion}[theorem]{Criterion}
\newtheorem{definition}[theorem]{Definition}
\newtheorem{example}[theorem]{Example}
\newtheorem{exercise}[theorem]{Exercise}
\newtheorem{lemma}[theorem]{Lemma}
\newtheorem{notation}[theorem]{Notation}
\newtheorem{problem}[theorem]{Problem}
\newtheorem{proposition}[theorem]{Proposition}
\newtheorem{remark}[theorem]{Remark}
\newtheorem{solution}[theorem]{Solution}
\newtheorem{summary}[theorem]{Summary}
\newenvironment{proof}[1][Proof]{\noindent\textbf{#1.} }{\ \rule{0.5em}{0.5em}}
\begin{document}
\preprint{HEP/123-qed}
\title[Short title for running header]{REV\TeX4}
\author{First Author}
\affiliation{My Institution}
\author{Second Author}
\affiliation{My Institution}
\author{Third Author}
\affiliation{Other Institution}
\keywords{one two three}
\pacs{PACS number}

\begin{abstract}
Shell document for REV\TeX{} 4.

\end{abstract}
\volumeyear{year}
\volumenumber{number}
\issuenumber{number}
\eid{identifier}
\date[Date text]{date}
\received[Received text]{date}

\revised[Revised text]{date}

\accepted[Accepted text]{date}

\published[Published text]{date}

\startpage{101}
\endpage{102}
\maketitle
\tableofcontents


\subsubsection{The Number Operator}

If a state describes on the average $N$ particles then%
\begin{equation}
f_{D}\left(  \mathcal{N}\right)  =\operatorname*{tr}\left\{  \left\langle
\mathcal{N}_{0}D_{0}\right\rangle _{0}\right\}  +\operatorname*{tr}\left\{
\left\langle \mathcal{N}_{2}D_{2}\right\rangle _{0}\right\}  =N,
\end{equation}
where the number operator $\mathcal{N}$ is
\begin{equation}
\mathcal{N}=\sum_{1\leq j\leq r}n_{j}=\mathcal{N}_{0}+\mathcal{N}%
_{2}=\mathcal{N}=\frac{r}{2}I_{\mathcal{F}}+i\sum_{1\leq j\leq r}x_{j}x_{j+r}.
\end{equation}
If it describes $N$ on the average particles then
\begin{equation}
f_{D}\left(  \mathcal{N}\right)  =\operatorname*{tr}\left\{  \left\langle
\mathcal{N}_{0}D_{0}\right\rangle _{0}\right\}  +\operatorname*{tr}\left\{
\left\langle \mathcal{N}_{2}D_{2}\right\rangle _{0}\right\}  =N,
\end{equation}
where the number operator $\mathcal{N}$ is
\begin{equation}
\mathcal{N}=\sum_{1\leq j\leq r}n_{j}=\mathcal{N}_{0}+\left\langle
\mathcal{N}\right\rangle _{2}=\mathcal{N}=\frac{r}{2}I_{\mathcal{F}}%
+i\sum_{1\leq j\leq r}x_{j}x_{j+r}.
\end{equation}
If a state $D\equiv f_{D}$ describes $N$ particles on the average then
\begin{align}
f_{D}\left(  \mathcal{N}\right)   &  =f_{D}\left(  \frac{r}{2}I_{\mathcal{F}%
}+i\sum_{1\leq j\leq r}x_{j}x_{j+r}\right)  =N\\
f_{D}\left(  \frac{r}{2}I_{\mathcal{F}}+i\sum_{1\leq j\leq r}x_{j}%
x_{j+r}\right)   &  =\frac{r}{2}+i\sum_{1\leq j\leq r}f_{D}\left(
x_{j}x_{j+r}\right) \\
\sum_{1\leq j\leq r}f_{D}\left(  x_{j}x_{j+r}\right)   &  =-i\left(
N-\frac{r}{2}\right)
\end{align}


If a state $D\equiv f_{D}$ describes $N$ particles then
\begin{gather}
\left[  \mathcal{N},D\right]  =0\,\&\,\operatorname*{tr}\left\{  \left\langle
\mathcal{N}D\right\rangle _{0}\right\}  =N\nonumber\\
\Updownarrow\nonumber\\
\operatorname*{tr}\left\{  \left\langle \mathcal{N}^{2}D\right\rangle
_{0}\right\}  -\left(  \operatorname*{tr}\left\{  \left\langle \mathcal{N}%
D\right\rangle _{0}\right\}  \right)  ^{2}=0,
\end{gather}
i.e.
\begin{equation}
f_{D}\left(  \mathcal{N}\right)  ^{2}-f_{D}\left(  \mathcal{N}^{2}\right)  =0.
\end{equation}


The number operator expressed in the $\left\{  x_{j}\right\}  $ bases involves
$r$, which is infinity when $r=\infty,$ however the back transformation works
as one gets $\infty-\infty$ in the form of $r-r$ as can be seen from%
\begin{align}
\mathcal{N}  &  =\frac{r}{2}I_{\mathcal{F}}+i\sum_{1\leq j\leq r}\frac
{1}{\sqrt{2}}\left(  a_{j}^{\dagger}+a_{j}\right)  \frac{i}{\sqrt{2}}\left(
a_{j}^{\dagger}-a_{j}\right) \nonumber\\
&  =\frac{r}{2}I_{\mathcal{F}}+i\sum_{1\leq j\leq r}\frac{i}{2}\left(
-a_{j}^{\dagger}a_{j}+a_{j}a_{j}^{\dagger}\right)  =\frac{r}{2}I_{\mathcal{F}%
}+\frac{1}{2}\sum_{1\leq j\leq r}\left[  a_{j}^{\dagger},a_{j}\right]
\nonumber\\
&  =\frac{r}{2}I_{\mathcal{F}}-\frac{1}{2}\sum_{1\leq j\leq r}\left(
-a_{j}^{\dagger}a_{j}+I_{\mathcal{F}}-a_{j}^{\dagger}a_{j}\right) \nonumber\\
&  =\frac{r}{2}I_{\mathcal{F}}-\sum_{1\leq j\leq r}\left(  \frac{1}%
{2}I_{\mathcal{F}}-a_{j}^{\dagger}a_{j}\right) \nonumber\\
&  =\frac{r}{2}I_{\mathcal{F}}-\frac{r}{2}I_{\mathcal{F}}+\sum_{1\leq j\leq
r}a_{j}^{\dagger}a_{j}\nonumber\\
&  =\mathcal{N}.
\end{align}
$\lceil$%
\begin{equation}
ix_{j}x_{j+r}=\frac{1}{2}I_{\mathcal{F}}-a_{j}^{\dagger}a_{j}=\frac{1}%
{2}\left[  a_{j}^{\dagger},a_{j}\right]
\end{equation}
$\rfloor$

The square of the number operator is given by
\begin{align}
\mathcal{N}^{2}  &  =\sum_{1\leq j,k\leq r}n_{j}n_{k}=\sum_{1\leq j\leq
r}n_{j}+2\sum_{1\leq j<k\leq r}n_{j}n_{k}=\sum_{1\leq j\leq r}n_{j}%
+2\sum_{1\leq j<k\leq r}a_{j}^{\dagger}a_{j}a_{k}^{\dagger}a_{k}\\
&  =\sum_{1\leq j\leq r}n_{j}+2\sum_{1\leq j<k\leq r}a_{j}^{\dagger}%
a_{k}^{\dagger}a_{k}a_{j}=\sum_{1\leq j\leq r}n_{j}+2\sum_{1\leq j<k\leq
r}n_{jk}=\mathcal{N}+2\mathcal{N}_{2},
\end{align}
which gives in the Clifford basis
\begin{align}
\mathcal{N}^{2}  &  =\left(  \frac{r}{2}I_{\mathcal{F}}+i\sum_{1\leq j\leq
r}x_{j}x_{j+r}\right)  \left(  \frac{r}{2}I_{\mathcal{F}}+i\sum_{1\leq k\leq
r}x_{k}x_{k+r}\right) \nonumber\\
&  =\frac{r^{2}}{4}I_{\mathcal{F}}+i\frac{r}{2}\sum_{1\leq j\leq r}%
x_{j}x_{j+r}+\sum_{1\leq j,k\leq r}x_{j}x_{k}x_{j+r}x_{k+r}\nonumber\\
&  =\left(  \frac{r^{2}}{4}+\frac{r}{4}\right)  I_{\mathcal{F}}+i\frac{r}%
{2}\sum_{1\leq j\leq r}x_{j}x_{j+r}+2\sum_{1\leq j<k\leq r}x_{j}x_{k}%
x_{j+r}x_{k+r}\nonumber\\
&  =\frac{r\left(  r+1\right)  }{4}I_{\mathcal{F}}+i\frac{r}{2}\sum_{1\leq
j\leq r}x_{j}x_{j+r}+2\sum_{1\leq j<k\leq r}x_{j}x_{k}x_{j+r}x_{k+r}.
\end{align}


Thus a state describes a fixed number of particles iff
\begin{equation}
f_{D}\left(  \mathcal{N}\right)  ^{2}-f_{D}\left(  \mathcal{N}^{2}\right)
=f_{D}\left(  \mathcal{N}\right)  ^{2}-f_{D}\left(  \mathcal{N}\right)
-2f_{D}\left(  \mathcal{N}_{2}\right)  =f_{D}\left(  \mathcal{N}\right)
\left(  f_{D}\left(  \mathcal{N}\right)  -1\right)  -2f_{D}\left(
\mathcal{N}_{2}\right)  =0,
\end{equation}
i.e.
\begin{equation}
f_{D}\left(  \mathcal{N}\right)  \left(  f_{D}\left(  \mathcal{N}\right)
-1\right)  =2f_{D}\left(  \mathcal{N}_{2}\right)  .
\end{equation}
In detail
\begin{align}
f_{D}\left(  \mathcal{N}\right)  ^{2}-f_{D}\left(  \mathcal{N}^{2}\right)   &
=f_{D}\left(  \frac{r}{2}I_{\mathcal{F}}+i\sum_{1\leq j\leq r}x_{j}%
x_{j+r}\right)  ^{2}-f_{D}\left(  \frac{r^{2}}{4}I_{\mathcal{F}}+i\sum_{1\leq
j\leq r}x_{j}x_{j+r}-\sum_{1\leq j,k\leq r}x_{j}x_{k}x_{j+r}x_{k+r}\right) \\
&  =\frac{r^{2}}{4}+ri\sum_{1\leq j\leq r}f_{D}\left(  x_{j}x_{j+r}\right)
-\left(  \sum_{1\leq j\leq r}f_{D}\left(  x_{j}x_{j+r}\right)  \right)
^{2}-\frac{r^{2}}{4}-ir\sum_{1\leq j\leq r}f_{D}\left(  x_{j}x_{j+r}\right)
+\sum_{1\leq j,k\leq r}f_{D}\left(  x_{j}x_{k}x_{j+r}x_{k+r}\right) \\
&  =\sum_{1\leq j,k\leq r}f_{D}\left(  x_{j}x_{k}x_{j+r}x_{k+r}\right)
-\left(  \sum_{1\leq j\leq r}f_{D}\left(  x_{j}x_{j+r}\right)  \right)  ^{2}%
\end{align}
Thus
\begin{align}
&  f_{D}\left(  \mathcal{N}\right)  ^{2}-f_{D}\left(  \mathcal{N}^{2}\right)
=0\\
&  \Rightarrow\sum_{1\leq j,k\leq r}f_{D}\left(  x_{j}x_{k}x_{j+r}%
x_{k+r}\right)  -\left(  \sum_{1\leq j\leq r}f_{D}\left(  x_{j}x_{j+r}\right)
\right)  ^{2}=0\\
&  \Rightarrow\sum_{1\leq j,k\leq r}f_{D}\left(  x_{j}x_{k}x_{j+r}%
x_{k+r}\right)  =\left(  \sum_{1\leq j\leq r}f_{D}\left(  x_{j}x_{j+r}\right)
\right)  ^{2}.
\end{align}
$\lceil$Checking this in another way gives%
\begin{align}
0  &  =f_{D}\left(  \mathcal{N}\right)  ^{2}-f_{D}\left(  \mathcal{N}%
^{2}\right)  =\left(  f_{D}\left(  \sum_{1\leq j\leq r}n_{j}\right)  \right)
^{2}-f_{D}\left(  \left(  \sum_{1\leq j\leq r}n_{j}\right)  ^{2}\right) \\
&  =\left(  \sum_{1\leq j\leq r}f_{D}\left(  n_{j}\right)  \right)  ^{2}%
-f_{D}\left(  \left(  \sum_{1\leq j,k\leq r}n_{j}n_{k}\right)  \right)
=\left(  \sum_{1\leq j\leq r}f_{D}\left(  n_{j}\right)  \right)  ^{2}%
-f_{D}\left(  \left(  2\sum_{1\leq j<k\leq r}n_{j}n_{k}+\sum_{1\leq j\leq
r}n_{j}\right)  \right) \\
&  =2\sum_{1\leq j<k\leq r}f_{D}\left(  n_{j}\right)  f_{D}\left(
n_{kj}\right)  +\sum_{1\leq j\leq r}\left(  f_{D}\left(  n_{j}\right)
\right)  ^{2}-2\sum_{1\leq j<k\leq r}f_{D}\left(  n_{j}n_{k}\right)
-\sum_{1\leq j\leq r}f_{D}\left(  n_{j}\right) \\
&  =\sum_{1\leq j\leq r}\left(  f_{D}\left(  n_{j}\right)  \right)  ^{2}%
+2\sum_{1\leq j<k\leq r}f_{D}\left(  n_{j}\right)  f_{D}\left(  n_{k}\right)
-\sum_{1\leq j\leq r}f_{D}\left(  n_{j}\right)  -2\sum_{1\leq j<k\leq r}%
f_{D}\left(  n_{j}n_{k}\right)
\end{align}%
\begin{align}
&  \sum_{1\leq j\leq r}\left(  f_{D}\left(  \left(  \frac{1}{2}I_{\mathcal{F}%
}+ix_{j}x_{j+r}\right)  \right)  \right)  ^{2}+2\sum_{1\leq j<k\leq r}%
f_{D}\left(  \left(  \frac{1}{2}I_{\mathcal{F}}+ix_{j}x_{j+r}\right)  \right)
f_{D}\left(  \left(  \frac{1}{2}I_{\mathcal{F}}+ix_{k}x_{k+r}\right)  \right)
\\
&  -\sum_{1\leq j\leq r}f_{D}\left(  \left(  \frac{1}{2}I_{\mathcal{F}}%
+ix_{j}x_{j+r}\right)  \right)  -2\sum_{1\leq j<k\leq r}f_{D}\left(  \left(
\frac{1}{2}I_{\mathcal{F}}+ix_{j}x_{j+r}\right)  \left(  \frac{1}%
{2}I_{\mathcal{F}}+ix_{k}x_{k+r}\right)  \right) \\
&  =\sum_{1\leq j\leq r}\left(  \frac{1}{2}+if_{D}\left(  x_{j}x_{j+r}\right)
\right)  ^{2}+2\sum_{1\leq j<k\leq r}\left(  \frac{1}{2}+if_{D}\left(
x_{j}x_{j+r}\right)  \right)  \left(  \frac{1}{2}+if_{D}\left(  x_{k}%
x_{k+r}\right)  \right) \\
&  -\sum_{1\leq j\leq r}\left(  \frac{1}{2}+if_{D}\left(  x_{j}x_{j+r}\right)
\right)  -2\sum_{1\leq j<k\leq r}f_{D}\left(  \frac{1}{4}+i\frac{1}{2}%
x_{j}x_{j+r}+i\frac{1}{2}x_{k}x_{k+r}-x_{j}x_{j+r}x_{k}x_{k+r}\right) \\
&  =\sum_{1\leq j\leq r}\left(  \frac{1}{4}+if_{D}\left(  x_{j}x_{j+r}\right)
-f_{D}\left(  x_{j}x_{j+r}\right)  ^{2}\right)  +2\sum_{1\leq j<k\leq
r}\left(  \frac{1}{4}+\frac{1}{2}if_{D}\left(  x_{j}x_{j+r}\right)  +\frac
{1}{2}if_{D}\left(  x_{k}x_{k+r}\right)  -f_{D}\left(  x_{j}x_{j+r}\right)
f_{D}\left(  x_{k}x_{k+r}\right)  \right) \\
&  -\sum_{1\leq j\leq r}\left(  \frac{1}{2}+if_{D}\left(  x_{j}x_{j+r}\right)
\right)  -2\sum_{1\leq j<k\leq r}\left(  \frac{1}{4}+f_{D}\left(  i\frac{1}%
{2}x_{j}x_{j+r}+i\frac{1}{2}x_{k}x_{k+r}-x_{j}x_{j+r}x_{k}x_{k+r}\right)
\right) \\
&  =\sum_{1\leq j\leq r}\left(  \frac{1}{4}+if_{D}\left(  x_{j}x_{j+r}\right)
-f_{D}\left(  x_{j}x_{j+r}\right)  ^{2}\right)  +2\sum_{1\leq j<k\leq
r}\left(  \frac{1}{2}if_{D}\left(  x_{j}x_{j+r}\right)  +\frac{1}{2}%
if_{D}\left(  x_{k}x_{k+r}\right)  -f_{D}\left(  x_{j}x_{j+r}\right)
f_{D}\left(  x_{k}x_{k+r}\right)  \right) \\
&  -\sum_{1\leq j\leq r}\left(  \frac{1}{2}+if_{D}\left(  x_{j}x_{j+r}\right)
\right)  -2\sum_{1\leq j<k\leq r}\left(  f_{D}\left(  i\frac{1}{2}x_{j}%
x_{j+r}+i\frac{1}{2}x_{k}x_{k+r}-x_{j}x_{j+r}x_{k}x_{k+r}\right)  \right) \\
&  =\sum_{1\leq j\leq r}\left(  \frac{1}{4}+if_{D}\left(  x_{j}x_{j+r}\right)
-f_{D}\left(  x_{j}x_{j+r}\right)  ^{2}\right)  +\sum_{1\leq j<k\leq r}\left(
if_{D}\left(  x_{j}x_{j+r}\right)  +if_{D}\left(  x_{k}x_{k+r}\right)
-2f_{D}\left(  x_{j}x_{j+r}\right)  f_{D}\left(  x_{k}x_{k+r}\right)  \right)
\\
&  -\sum_{1\leq j\leq r}\left(  \frac{1}{2}+if_{D}\left(  x_{j}x_{j+r}\right)
\right)  -\sum_{1\leq j<k\leq r}\left(  if_{D}\left(  x_{j}x_{j+r}\right)
+if_{D}\left(  x_{k}x_{k+r}\right)  -2f_{D}\left(  x_{j}x_{j+r}x_{k}%
x_{k+r}\right)  \right) \\
&  =-\sum_{1\leq j\leq r}\left(  \frac{1}{4}+f_{D}\left(  x_{j}x_{j+r}\right)
^{2}\right)  +2\sum_{1\leq j<k\leq r}\left(  f_{D}\left(  x_{j}x_{j+r}%
x_{k}x_{k+r}\right)  -f_{D}\left(  x_{j}x_{j+r}\right)  f_{D}\left(
x_{k}x_{k+r}\right)  \right) \\
&  =-\frac{r}{4}-\sum_{1\leq j\leq r}f_{D}\left(  x_{j}x_{j+r}\right)
^{2}+\sum_{1\leq j,k\leq r}\left(  f_{D}\left(  x_{j}x_{j+r}x_{k}%
x_{k+r}\right)  -f_{D}\left(  x_{j}x_{j+r}\right)  f_{D}\left(  x_{k}%
x_{k+r}\right)  \right)  -\sum_{1\leq j\leq r}\left(  -\frac{1}{4}%
-f_{D}\left(  x_{j}x_{j+r}\right)  ^{2}\right) \\
&  =\sum_{1\leq j,k\leq r}\left(  f_{D}\left(  x_{j}x_{j+r}x_{k}%
x_{k+r}\right)  -f_{D}\left(  x_{j}x_{j+r}\right)  f_{D}\left(  x_{k}%
x_{k+r}\right)  \right)
\end{align}
$\rfloor$

The pair operator $\mathcal{N}_{2}$ can be expanded as
\begin{align}
\mathcal{N}_{2}  &  =\sum_{1\leq j<k\leq r}n_{j}n_{k}=\sum_{1\leq j<k\leq
r}\left(  \frac{1}{2}I_{\mathcal{F}}+ix_{j}x_{j+r}\right)  \left(  \frac{1}%
{2}I_{\mathcal{F}}+ix_{k}x_{k+r}\right) \\
&  =\frac{1}{4}\binom{r}{2}I_{\mathcal{F}}+2i\sum_{1\leq j\leq r}x_{j}%
x_{j+r}-\sum_{1\leq j<k\leq r}x_{j}x_{j+r}x_{k}x_{k+r}=\left\langle
\mathcal{N}_{2}\right\rangle _{0}+\left\langle \mathcal{N}_{2}\right\rangle
_{2}+\left\langle \mathcal{N}_{2}\right\rangle _{4}.
\end{align}


The number operator is an unbounded positive operator that is densly defined
and has a spectral resolution
\begin{equation}
\mathcal{N}=\sum_{0\leq N\leq\infty}NP_{N},
\end{equation}
where
\begin{align}
P_{N}  &  =\left(  P_{N}\right)  ^{\dagger};P_{N}=\left(  P_{N}\right)  ^{2}\\
P_{N}  &  =\sum_{0\leq j\leq\infty}p_{j};\,p_{j}\text{ are atomic
projectors}\\
\operatorname*{tr}\left\{  P_{N}\right\}   &  =\binom{r}{N}%
\end{align}
and thus are infinite dimensional if $r=\infty$ and $N>0.$ One can define the
phase operators
\begin{equation}
\Phi\left(  s\right)  =e^{i\mathcal{N}s}%
\end{equation}
that form a one parameter unitary group. One can define functions $h\left(
\mathcal{N}\right)  $ of $\mathcal{N}$ by
\begin{equation}
h\left(  \mathcal{N}\right)  =\sum_{0\leq N\leq\infty}h\left(  N\right)  P_{N}%
\end{equation}
so that
\begin{equation}
h\left(  \mathcal{N}\right)  ^{2}=\sum_{0\leq N\leq\infty}h\left(  N\right)
^{2}P_{N}%
\end{equation}
and if $f_{D}$ is a state describing a fixed number of particles
\begin{equation}
f_{D}\left(  h\left(  \mathcal{N}\right)  ^{2}\right)  =f_{D}\left(  h\left(
\mathcal{N}\right)  \right)  ^{2}.
\end{equation}
$\lceil$If $f_{D}$ is normal state then
\begin{equation}
f_{D}\left(  X\right)  =\operatorname*{tr}\left\{  XD\right\}
\end{equation}
and if $f_{D}$ is a state describing a fixed number of particles $N$ then
\begin{align}
D  &  =\sum_{1\leq j\leq\infty}d_{j}P_{D_{j}}\\
P_{N}P_{D_{j}}  &  =P_{D_{j}}P_{N}=P_{D_{j}};1\leq j\leq\infty\\
P_{M}P_{D_{j}}  &  =P_{D_{j}}P_{M}=0;1\leq j\leq\infty,M\neq N.
\end{align}
Thus
\begin{align}
f_{D}\left(  h\left(  \mathcal{N}\right)  \right)   &  =\operatorname*{tr}%
\left\{  h\left(  \mathcal{N}\right)  D\right\}  =\operatorname*{tr}\left\{
\sum_{0\leq M\leq\infty}\sum_{1\leq j\leq\infty}h\left(  M\right)
d_{j}P_{D_{j}}\right\}  =\operatorname*{tr}\left\{  \sum_{1\leq j\leq\infty
}h\left(  N\right)  d_{j}P_{D_{j}}\right\}  =\sum_{1\leq j\leq\infty}h\left(
N\right)  d_{j}=h\left(  N\right) \\
f_{D}\left(  h\left(  \mathcal{N}\right)  ^{2}\right)   &  =\operatorname*{tr}%
\left\{  h\left(  \mathcal{N}\right)  ^{2}D\right\}  =\operatorname*{tr}%
\left\{  \sum_{0\leq M\leq\infty}\sum_{1\leq j\leq\infty}h\left(  M\right)
^{2}d_{j}P_{D_{j}}\right\}  =\operatorname*{tr}\left\{  \sum_{1\leq
j\leq\infty}h\left(  N\right)  ^{2}d_{j}P_{D_{j}}\right\}  =\sum_{1\leq
j\leq\infty}h\left(  N\right)  ^{2}d_{j}=h\left(  N\right)  ^{2}.
\end{align}
Want $h\left(  \mathcal{N}\right)  $ so that
\begin{align}
h\left(  \mathcal{N}\right)   &  =\lim\limits_{r\rightarrow\infty}h\left(
\frac{r}{2}I_{\mathcal{F}}+i\sum_{1\leq j\leq r}x_{j}x_{j+r}\right) \\
\operatorname*{tr}\left\{  h\left(  \mathcal{N}\right)  \right\}   &
=\operatorname*{tr}\left\{  h\left(  \lim\limits_{r\rightarrow\infty}h\left(
\frac{r}{2}I_{\mathcal{F}}+i\sum_{1\leq j\leq r}x_{j}x_{j+r}\right)  \right)
\right\}  <\infty.
\end{align}
%

\begin{equation}
\mathcal{N}^{2}=\left(  \sum_{0\leq N\leq\infty}NP_{N}\right)  ^{2}%
=\sum_{0\leq N\leq\infty}N^{2}P_{N}%
\end{equation}%
\begin{align}
e^{\mathcal{N}}  &  =\sum_{0\leq N\leq\infty}P_{N}+\sum_{0\leq N\leq\infty
}NP_{N}+\frac{1}{2}\sum_{0\leq N\leq\infty}N^{2}P_{N}+\cdots\\
&  =e^{N}I_{\mathcal{F}}%
\end{align}


\section{Constrained Hamiltonian}

If the stationary states of the Hamiltonian are constrained to have on the
average to describe $N$ electrons then the Lagrangian describing this
constrained optimization is given by
\begin{align}
\mathcal{L}\left(  H\boldsymbol{,}\mathcal{N},\lambda\right)   &
=H-\lambda\mathcal{N}=\eta_{0}I_{\mathcal{F}}+i\sum_{1\leq j_{1}<j_{2}\leq
2r}\eta_{2j_{1}j_{2}}x_{j_{1}}x_{j_{2}}+\sum_{1\leq j_{1}<j_{2}<j_{3}%
<j_{4}\leq2r}\eta_{4j_{1}j_{2}j_{3}j_{4}}x_{j_{1}}x_{j_{2}}x_{j_{3}}x_{j_{4}%
}-\lambda\left\{  \left(  \frac{r}{2}I_{\mathcal{F}}+i\sum_{1\leq j\leq
r}x_{j}x_{j+r}\right)  \right\} \\
&  =\left\{  h_{0}+\frac{1}{2}\left(  \operatorname*{Tr}h_{1}\right)
+\frac{1}{4}\operatorname*{Tr}\left\{  h_{2}\right\}  -\frac{\lambda r}%
{2}\right\}  I_{\mathcal{F}}+i\sum_{1\leq j_{1}<j_{2}\leq2r}\eta_{2j_{1}j_{2}%
}x_{j_{1}}x_{j_{2}}-\lambda i\sum_{1\leq j\leq r}x_{j}x_{j+r}+\sum_{1\leq
j_{1}<j_{2}<j_{3}<j_{4}\leq2r}\eta_{4j_{1}j_{2}j_{3}j_{4}}x_{j_{1}}x_{j_{2}%
}x_{j_{3}}x_{j_{4}}%
\end{align}


\subsection{Time Evolution}

\qquad%
\begin{align}
U\left(  t\right)   &  =\exp-i\left\{  \left(  \eta_{0}I_{\mathcal{F}}%
+i\eta_{2}+\eta_{4}\right)  t\right\}  =e^{-i\eta_{0}t}\exp i\left\{  \left(
i\eta_{2}+\eta_{4}\right)  t\right\}  \in U\left(  1\right)  \times U\left(
\mathcal{H}\right)  /U\left(  1\right)  =U\left(  1\right)  \times SU\left(
\mathcal{H}\right) \\
D\left(  t\right)   &  =U\left(  t\right)  D_{0}U\left(  t\right)  ^{\dagger
}=\exp-i\left\{  \left(  i\eta_{2}+\eta_{4}\right)  t\right\}  D_{0}\exp
i\left\{  \left(  -i\eta_{2}^{\dagger}+\eta_{4}\right)  t\right\} \\
&  =\exp\left\{  \left(  \eta_{2}-i\eta_{4}\right)  t\right\}  D_{0}%
\exp\left\{  \left(  \eta_{2}+i\eta_{4}\right)  t\right\}  =\left(
\exp-i\Lambda t\right)  D_{0}\left(  \exp i\Lambda t\right) \\
\Lambda &  =\left(  i\eta_{2}+\eta_{4}\right)  =\left(  i\eta_{2}+\eta
_{4}\right)  ^{\dagger}=-i\eta_{2}^{t}+\eta_{4}^{t}=i\eta_{2}+\eta_{4}%
\end{align}%
\begin{align}
&  \hat{U}\left(  t\right)  \left(  D_{0}\right)  =\left(  \exp-i\Lambda
t\right)  D_{0}\left(  \exp i\Lambda t\right)  =\left(  \exp-i\hat{\Lambda
}t\right)  \left(  D_{0}\right) \\
&  =\left(  I_{\mathcal{F}}-i\Lambda t-\frac{1}{2}\Lambda^{2}t^{2}%
+\cdots\right)  D_{0}\left(  I_{\mathcal{F}}+i\Lambda t-\frac{1}{2}\Lambda
^{2}t^{2}+\cdots\right) \\
&  =D_{0}-i\Lambda D_{0}t+iD_{0}\Lambda t-\frac{1}{2}\left(  \Lambda^{2}%
D_{0}+D_{0}\Lambda^{2}\right)  t^{2}+\cdots
\end{align}%
\begin{align}
&  \frac{d}{dt}\left\{  \left(  \exp-i\Lambda t\right)  D_{0}\left(  \exp
i\Lambda t\right)  \right\} \\
&  =-i\Lambda D_{0}+iD_{0}\Lambda-\left(  \Lambda^{2}D_{0}+D_{0}\Lambda
^{2}\right)  t+\cdots\\
&  =-i\Lambda\left\{  D_{0}+i\Lambda D_{0}t+\cdots\right\}  +\left\{
D_{0}+i\Lambda D_{0}t+\cdots\right\}  i\Lambda\\
&  =-i\Lambda D\left(  t\right)  +iD\left(  t\right)  \Lambda=i\left[
D\left(  t\right)  ,\Lambda\right]
\end{align}
i.e.
\begin{equation}
-i\frac{dD\left(  t\right)  }{dt}=\left[  D\left(  t\right)  ,\Lambda\right]
.
\end{equation}
Stationary states are thus characterized by
\begin{equation}
\left[  D\left(  t\right)  ,\Lambda\right]  =0
\end{equation}
If $D_{j}$ is a vector state (extreme) then
\begin{equation}
D_{j}=\left\vert \Psi_{j}\right\rangle \left\langle \Psi_{j}\right\vert
\end{equation}
and%
\begin{align}
\left[  \left\vert \Psi_{j}\right\rangle \left\langle \Psi_{j}\right\vert
,\Lambda\right]   &  =0=\left\vert \Psi_{j}\right\rangle \left\langle \Psi
_{j}\right\vert \Lambda-\Lambda\left\vert \Psi_{j}\right\rangle \left\langle
\Psi_{j}\right\vert =\left\vert \Psi_{j}\right\rangle \left\langle \Lambda
\Psi_{j}\right\vert -\left\vert \Lambda\Psi_{j}\right\rangle \left\langle
\Psi_{j}\right\vert \\
&  \Rightarrow\Lambda\left\vert \Psi_{j}\right\rangle =\lambda_{j}\left\vert
\Psi_{j}\right\rangle .
\end{align}
If $\Lambda$ only has a discrete spectra
\begin{equation}
\sum_{1\leq j\leq\infty}\lambda_{j}=0.
\end{equation}


The expectation values vary in time as
\begin{equation}
f_{D\left(  t\right)  }\left(  A\right)  =\operatorname{Tr}\left\{  AD\left(
t\right)  \right\}  =\operatorname{Tr}\left\{  A\left(  \exp-i\hat{\Lambda
}t\right)  \left(  D_{0}\right)  \right\}  =\operatorname{Tr}\left\{  \left(
\exp i\hat{\Lambda}t\right)  \left(  A\right)  D_{0}\right\}
\end{equation}


The majorana number operator
\begin{equation}
\left\langle \mathcal{N}\right\rangle _{2}=i\sum_{1\leq j\leq r}x_{j}%
x_{j+r}=\frac{1}{2}\sum_{1\leq j\leq r}\left[  a_{j}^{\dagger},a_{j}\right]
\end{equation}
produces the one parameter group of transformations
\begin{equation}
\hat{U}_{\mathcal{N}}\left(  t\right)  \left(  D_{0}\right)  =\left(
\exp-i\left\langle \mathcal{N}\right\rangle _{2}t\right)  D_{0}\left(  \exp
i\left\langle \mathcal{N}\right\rangle _{2}t\right)  =\left(  \exp
-\widehat{i\left\langle \mathcal{N}\right\rangle _{2}}t\right)  \left(
D_{0}\right)
\end{equation}
where
\begin{equation}
\exp\left\{  -i\left\langle \mathcal{N}\right\rangle _{2}\right\}
=\exp\left\{  \sum_{1\leq j\leq r}x_{j}x_{j+r}\right\}  =\exp\left\{
J\right\}
\end{equation}
and $J$ is a complex unit.%
\begin{equation}
\exp\left\{  i\alpha\hat{A}\right\}  =\cos\left(  \alpha\right)  I+\hat{A}%
\sin\left(  \alpha\right)  ,
\end{equation}%
\begin{equation}
\exp\left\{  i\frac{\pi}{2}\sum_{1\leq j\leq r}x_{j}x_{j+r}\right\}
=\sum_{1\leq j\leq r}x_{j}x_{j+r},
\end{equation}


\section{Operators and Von Neumann Algebras}

The \emph{Resolvent set}, $\mathcal{R}\left(  A\right)  $\emph{ }of an
operator $A\in\mathcal{A}$ (c$^{\ast}$-algebra) is defined to be
\begin{equation}
\mathcal{R}\left(  A\right)  =\left\{  \left.  z\right\vert z\in
\mathbb{C},\left(  zI-A\right)  ^{-1}\,\exists\right\}
\end{equation}
and the $\emph{Spectrum}$, $Sp\left(  A\right)  $ of an operator
$A\in\mathcal{A}$ is the complement of $\mathcal{R}\left(  A\right)  $ i.e.
\begin{equation}
Sp\left(  A\right)  \equiv\sigma\left(  A\right)  =\mathcal{R}\left(
A\right)  ^{c}%
\end{equation}


\begin{theorem}
Spectral Theorem

For any normal element $A\in\mathcal{A},$ every repreentation of $\mathcal{A}$
is equivalent to a representation
\begin{equation}
\mathcal{H\,\,}=\mathcal{\,}\bigoplus\mathcal{H}_{j}%
\end{equation}
for which
\begin{equation}
\mathcal{H}_{j}=L^{2}\left(  \sigma\left(  A\right)  ,d\mu_{j}\right)
\end{equation}
and
\begin{equation}
\Pi\left(  A\right)  \varphi\left(  \alpha\right)  =\alpha\varphi\left(
\alpha\right)  .
\end{equation}

\end{theorem}

If $A=A^{\dagger}$ then
\begin{equation}
A=\int_{-\infty}^{\infty}\alpha dP_{A}\left(  \alpha\right)  =\int_{\min
Sp\left(  A\right)  }^{\max Sp\left(  A\right)  }\alpha dP_{A}\left(
\alpha\right)  =\int_{\min Sp\left(  A\right)  }^{\max Sp\left(  A\right)
}\alpha\rho_{A}\left(  \alpha\right)  d\alpha;\alpha\in\mathbb{R}%
\end{equation}
and
\begin{equation}
\Pi_{L^{2}\left(  \sigma\left(  A\right)  ,d\mu_{A}\right)  }\left(  A\right)
=\int_{\min Sp\left(  A\right)  }^{\max Sp\left(  A\right)  }\alpha d\mu
_{sA}\left(  \alpha\right)  =\int_{\min Sp\left(  A\right)  }^{\max Sp\left(
A\right)  }\alpha d\mu_{pA}\left(  \alpha\right)  +\int_{\min Sp\left(
A\right)  }^{\max Sp\left(  A\right)  }\alpha d\mu_{acA}\left(  \alpha\right)
+\int_{\min Sp\left(  A\right)  }^{\max Sp\left(  A\right)  }\alpha d\mu
_{sA}\left(  \alpha\right)  ,
\end{equation}
where%
\begin{align}
d\mu_{pA}\left(  \alpha\right)   &  =\left(  \sum_{1\leq m\leq\infty}%
a_{m}\delta\left(  \alpha-a_{m}\right)  \right)  d\alpha\\
d\mu_{acA}\left(  \alpha\right)   &  =f_{A}\left(  \alpha\right)
d\alpha;f_{A}\left(  \alpha\right)  \geq0\\
d\mu_{sA}\left(  \alpha\right)   &  =d\mu_{A}\left(  \alpha\right)  -d\mu
_{pA}\left(  \alpha\right)  -d\mu_{acA}\left(  \alpha\right) \\
\mu_{sA}\left(  \Delta\right)   &  \neq0\Leftrightarrow\mu_{pA}\left(
\Delta\right)  =0\,\&\,\mu_{acA}\left(  \Delta\right)  =0\\
\mu_{pA}\left(  \Delta\right)   &  \neq0\Leftrightarrow\mu_{sA}\left(
\Delta\right)  =0\,\&\,\mu_{acA}\left(  \Delta\right)  =0\\
\mu_{acA}\left(  \Delta\right)   &  \neq0\Leftrightarrow\mu_{pA}\left(
\Delta\right)  =0\,\&\,\mu_{sA}\left(  \Delta\right)  =0.
\end{align}
i.e. the measures $\mu_{acA},\mu_{pA}$ and $\mu_{sA}$ are orthogonal to each
other.
\begin{equation}
L^{2}\left(  \sigma\left(  A\right)  ,d\mu_{A}\right)  =L^{2}\left(
\sigma\left(  A\right)  ,d\mu_{pA}\right)  \oplus L^{2}\left(  \sigma\left(
A\right)  ,d\mu_{acA}\right)  \oplus L^{2}\left(  \sigma\left(  A\right)
,d\mu_{sA}\right)
\end{equation}
In general for any representation $\Pi_{\mathcal{H}}$
\begin{equation}
\mathcal{H}=\mathcal{H}_{pA}\mathcal{\oplus H}_{acA}\mathcal{\oplus H}_{sa}.
\end{equation}
so that
\begin{equation}
\Pi_{\mathcal{H}}\left(  A\right)  =\sum_{0\leq j\leq\infty}\alpha
_{j}\left\vert \Psi_{Aj}\right\rangle \left\langle \Psi_{Aj}\right\vert
+\int_{\min Sp\left(  A\right)  }^{\max Sp\left(  A\right)  }\alpha\rho
_{acA}\left(  \alpha\right)  d\alpha+\int_{\min Sp\left(  A\right)  }^{\max
Sp\left(  A\right)  }%
\end{equation}

\end{document}